\section{의존성에 따른 다음 작업과 참여 지점}
\label{sec:roadmap}

\subsection{현재 단일 우선순위: security ExpandVecA 의미론 수리}

최신 post-freeze gate는 \CurrentGate 다. 첫 의무는
\begin{quote}
\CurrentPriority
\end{quote}
이다. post-freeze NTT/Rq/HAETAE 운송과 actual \identifier{avec}의 paper \(a\)
해석은 컴파일되었다. 그러나 기존 security-side object는 actual SHAKE128/rejection
expansion도 paper \(qj\)도 아닌 synthetic generator다. 이 불일치를 고치지 않은
채 전체 KeyGen 식을 합성하면 concrete coefficient 44985와 401을 같다고 놓게 된다.
한편 frozen paper의 \FrozenPaperGate 는 별도 상태다. 순수 2-by-4 행곱
specialization이 post-freeze에 생겼어도 actual Verify helper loop와 CRT/freeze
절차 정리는 여전히 열려 있다.

\begin{figure}[ht]
\centering
\begin{tikzpicture}[node distance=6mm]
  \node[provennode] (snapshot) {1. actual matrix/finalizer\\snapshot};
  \node[provennode, below=of snapshot] (ntt) {2. full-NTT와 native \(R_q\)\\row product};
  \node[provennode, below=of ntt] (repr) {3. \(R_q\)--HAETAE\\dot representation};
  \node[provennode, below=of repr] (avec) {4. actual \identifier{avec}\(\to\)paper \(a\)\\paper \(qj\) 별도 합성};
  \node[provennode, below=of avec] (faithful) {5. actual-carrier structured product\\\(\equiv qj\pmod{2q}\)};
  \node[blockednode, below=of faithful] (security) {6. security ExpandVecA\\semantic equivalence (STOP)};
  \node[blockednode, below=of security] (paper) {7. security-model\\KG-1--KG-4 lift};
  \draw[flowarrow] (snapshot) -- (ntt);
  \draw[flowarrow] (ntt) -- (repr);
  \draw[flowarrow] (repr) -- (avec);
  \draw[flowarrow] (avec) -- (faithful);
  \draw[openarrow] (faithful) -- (security);
  \draw[openarrow] (security) -- (paper);
\end{tikzpicture}
\caption{post-freeze KeyGen의 현재 절단면. actual-carrier 전용 coefficient-product까지
닫혔지만, generic security ExpandVecA/\identifier{matrix_vec_mul} 의미론에서 중단했다.}
\label{fig:next-order}
\end{figure}

codec 선에서는 Week~15에 fixed all-6 반환 실행의 failure exclusion까지 닫혔다.
그러나 fixed-input termination/losslessness는 열려 있고, \(h\) suffix codec은
이번 one-week MINCORE 범위 전환으로 \statusdeferred 되었다. 이 두 사실은 현재
KeyGen blocker와 독립적이다.

\subsection{1단계 완료: encoder}

\claimid{OBL-RANS-ENCODE-REFINEMENT}는 실제
\identifier{_rans_encode}를 여는
\theoremname{actual_rans_encode_trace_closure}와 그 확장된 따름정리(corollary)
\theoremname{actual_rans_encode_trace_refinement}로 닫혔다. 성공 분기(success
branch)의 사후조건(postcondition)은 다음을 동시에 준다.

\begin{itemize}
  \item actual output suffix가 canonical symbol list의
        \identifier{trace_bytes}와 점별로 같음;
  \item 바깥 반복(outer iteration)마다 actual word state가 suffix의 순수
        \identifier{encode_trace} state와 같음;
  \item actual \(off\)와 trace byte 길이가
        \(off+|\operatorname{TraceBytes}|=1024\)를 만족함;
  \item 이미 증명된 final 4-byte state 직렬화(serialization)를 normalization
        suffix 앞에 합성하여 segment 순서를 확정함;
  \item 쓰지 않은 접두부(unwritten prefix)의 frame을 유지함;
  \item 실패 branch를 숨기지 않음.
\end{itemize}

Week~10의 array/list suffix 점화식(recurrence), actual reciprocal word step,
내부 byte/frame/no-underflow, final serialization leaf와 성공 offset/size 범위는
Week~11의 tail-aware 바깥/안쪽 불변식(outer/inner invariant), generated table
load--protect--word-update 보조정리(lemma), final-store 합성(composition)에
연결되었다. 따라서 임의의 valid trace나 새 증인 버퍼(witness buffer)가 아니라
실제 반환 버퍼의 suffix가 정준 trace 바이트(canonical trace bytes)와 일치한다.

다만 이는 Hoare 부분정당성(Hoare partial correctness)이다. 실제 실행이 성공
분기(success branch)에 도달한다는 성공 증인(success witness)과 종료성(termination)은
이 Week~11 정리 하나에서 따라오지 않는다. Week~15는 fixed all-6 반환 실행의
failure를 별도 정리로 배제했지만 종료성은 여전히 남는다.

\subsection{2단계 완료: decoder}

\claimid{OBL-RANS-DECODE-REFINEMENT}는 실제
\identifier{_rans_decode}를 여는
\theoremname{actual_rans_decode_trace_refinement}로 닫혔다. exact-trace 입력
관계(input relation) 아래 다음이 귀납법(induction)으로 증명되었다.

\begin{itemize}
  \item iteration \(i\)에서 decoded prefix가 원래 symbol prefix와 같음;
  \item state/cursor가 pure \(x_i,c_i\)와 같음;
  \item actual table lookup이 구간 분할(interval partition)의 유일한 \(s_i\)를 선택함;
  \item nested normalization loop가 \(\beta_i\)를 정확한 순서로 소비함;
  \item 내부 loop 종료에서 \(state=2^{23}\)을 유도하고, 반환값에서
        \(off=size\), \(bad=0\)을 결론냄;
  \item decoded tail frame.
\end{itemize}

encoder와 decoder를 직접 lockstep으로 묶는 방법은 reverse/forward 방향과
중간 buffer copy 때문에 proof surface가 커진다. 이미 구성된 순수 trace를
공통 매개체로 삼아 두 refinement를 독립적으로 검증했다. 단, 이 decoder
정리(theorem)는 입력 trace를 전제로 받는 Hoare 부분정당성(Hoare partial
correctness)이며 종료성(termination)을 주지 않는다.

\subsection{3단계 완료: actual harness inverse}

\claimid{OBL-RANS-CORE-INVERSE}는
\sourcepath{Mode2RansCoreActualInverse.ec}의
\theoremname{actual_rans_encode_copy_decode_inverse}로 닫혔다. 이 정리(theorem)는
\sourcepath{Mode2RansActualInverse.ec}에 정의된 actual harness를 직접 대상으로
삼아 다음 의미론적 결론을 낸다. 두 파일명은 각각 ``harness 정의/control''과
``Week~13 semantic inverse''를 가리키며, 후자는 전자의 control-only 정리(theorem)를
의미론적 결과로 잘못 읽지 않도록 별도 증명된 것이다.

\begin{itemize}
  \item encoder 실패 분기(failure branch)를 숨기지 않고 decoder를 실행하지 않음;
  \item encoder 성공 suffix의 \identifier{segment_matches}를 actual copy의
        \identifier{slice_eq}에 합성함;
  \item copy된 buffer에서 decoder의 pointwise exact-read 전제를 산출함;
  \item actual decoder 반환에서 \(bad=0\), exact consumption, 1024-symbol
        등식(equality), tail frame을 얻음;
  \item 임의의 trace, 새 증인 버퍼(witness buffer), encoder 성공을 외부 전제로
        넣지 않음.
\end{itemize}

이 정리(theorem)는 성공 조건부 부분정확성(success-conditioned partial
correctness)이다. canonical all-zero array는 전제의 비공허성(non-vacuity)을
보이지만, encoder 성공의 존재나 전체 실행의 종료성(termination)을 자동으로
주지 않는다. Week~13 당시 성공 도달성은
\claimid{OBL-RANS-ACTUAL-SUCCESS-WITNESS}의 별도 의무였고, Week~15는 그중
fixed all-6 Hoare 사후조건만 닫았다.

\subsection{4단계 완료: production full-HBZ wrapper}

\claimid{OBL-SIG-HBZ-ENCODE-DECODE}는
\theoremname{signature_pack_unpack_hbz_full_actual_exact}와
\theoremname{signature_pack_unpack_hbz_full_inverse_mode2}로 닫혔다. actual
\identifier{_encode_hb_z1_full}과 \identifier{_decode_hb_z1_full}을 직접 여는
사슬은 다음을 결론낸다.

\begin{itemize}
  \item canonical HBZ coefficient를 actual prepare가 정준 기호열(canonical
        symbol stream)로 옮김;
  \item actual rANS core 정리(theorem)의 실패/성공 분기를 그대로 보존함;
  \item 성공 시 actual apply가 복원된 symbol을 원래 coefficient로 되돌림;
  \item wrapper-level size, state, 입력/출력 tail frame을 연결함;
  \item actual success witness나 decoder 결과를 외부 전제로 다시 요구하지 않음.
\end{itemize}

이 단계는 성공 조건부 부분정확성(success-conditioned partial correctness)으로
\statusproved 다 \cite{week14-report,hbz-full-wrapper-composition}. 실제 성공
가지의 일반 도달가능성(reachability)과 종료성(termination)은 별도
정리(theorem)다. Week~15는 fixed all-zero 입력의 반환 실행에 한해 failure를
배제한다.

\subsection{5단계 완료: fixed all-6 success 사후조건}

Week~15의 \theoremname{actual_rans_encode_all_six_success}는 actual failure가
trace length \(>1020\)을 요구한다는 정리와 all-6 trace의 \(\le1020\) budget을
모순시켜 반환하는 fixed-input 실행의 \(bad=0\)을 얻는다. focused/production
lift는 \(4\le size\le1024\), exact trace/frame과 zero coefficient round trip을
동시에 보인다 \cite{week15-report,rans-actual-success-witness}.

이 단계의 완료 판정은 fixed-input Hoare 부분정확성이다. termination,
losslessness, non-vacuous execution 또는 \identifier{phoare}로 승격하지 않는다.
원래 뒤따를 예정이던 \(h\) suffix, metadata/padding, full signature inverse와
Sign output tail reach는 Week~16 MINCORE 전환으로 각각 deferred 또는 blocked
상태를 유지한다.

\subsection{6단계 frozen 종료: KG--NTT--MUL 감사}

\theoremname{output_row_from_mode2_ntt_words}와
\theoremname{actual_m23_matrix_snapshot_rows_explicit}는
\identifier{KeygenM23ArithmeticSpec.output_row}를 actual full-invNTT/pointwise
word 표현으로 rewrite하고 두 active row에 적용하는 마지막 sound edge를 닫았다
\cite{week16-kg-ntt-mul-report}. frozen paper 시점에는 다음 두 화살표가 없었다.
\[
  \operatorname{array256\_mont}
    (\operatorname{full\_invntt}(\widehat a\odot\operatorname{full\_ntt}(p)))
  \;\not\Longrightarrow\;
  (\operatorname{full\_invntt}\widehat a)\star_{R_q}p,
\]
그리고 native \identifier{Rq.poly} 곱을 보안 모형의 integer-list
\identifier{poly_mul}/\identifier{matrix_vec_mul}로 옮기는 어댑터도 없다. 첫
화살표에는 odd-root orthogonality와 full-NTT convolution이 필요했다. 시간 제한
감사는 이를 새 전제로 가정하지 않고 \identifier{STOP-KG-NTT}로 종료했다.
post-freeze에 이 순수 대수와 \(R_q\)--HAETAE 표현 간극은 닫혔지만, 이 역사적
frozen 판정 자체는 변경하지 않는다.

\subsection{7단계 동결: accepted Sign core}

다음 세 actual helper의 순차 실행은 투명한 harness에서 합성되어 control
정리까지 컴파일됐다.
\begin{center}
\identifier{_sf_round_challenge_mode2}\quad
\identifier{_sf_z_check}\quad
\identifier{_sf_hint_mode2}
\end{center}
컴파일된 control theorem은
\sourcepath{Mode2SignAcceptedCore.ec}의
\begin{center}
\theoremname{actual_sign_accepted_core_branch_control_mode2}
\end{center}
이다. 그러나 headline S-1--S-7 정리는
\identifier{sf_challenge_mode2_highbits_lsb_sampleinball_correct} 부재로
\identifier{STOP-SIGN-CHAL-MODE2}에서 동결되었다.
반환하고 accept한 실행에 대해 commitment/challenge, response
\(z=y+(-1)^bcs\), 두 exact integer norm predicate와 coefficientwise hint 식을
도출해야 한다 \cite{week16-mincore-plan,claim-ledger}.

headline 정리는 아직 컴파일된 증명으로 세어서는 안 된다.
사전조건이 \identifier{reject=0}이나 최종 식을 가정해서도 안 된다. sampler
분포, signing-loop 종료성, packing과 전체 1474-byte signature는 이 단계의 범위
밖이다.

\subsection{8단계 frozen 종료: actual Verify helper chain}

\sourcepath{Mode2VerifyCoreSequence.ec}의
\theoremname{actual_verify_core_sequence_branch_control_mode2}와
\theoremname{actual_verify_core_sequence_input_snapshots_mode2}는 actual
Verify의 다섯 helper 호출 순서와 입력 snapshot을 고정한다.
\sourcepath{Mode2VerifyPrepareNorm.ec},
\sourcepath{Mode2VerifyRecover.ec},
\sourcepath{Mode2VerifyTailChallenge.ec}는 각각
\theoremname{verify_prepare_z1_wprime_mode2_word_exact},
\theoremname{actual_sign_verify_recover_w_z2_mode2_word_semantics},
\theoremname{verify_tail_exact_trace_mode2}와
\theoremname{poly_mismatch_mode2_word_exact}를 제공한다
\cite{week16-verify-report}. 따라서 V-1/V-2/V-5/V-6 projection, W64 norm gate,
tail trace/mismatch expression은 helper-local 범위에서 닫혔다.

그러나 actual 2-by-4
\identifier{_polyvec_ntt}/\identifier{_polymat_pointwise_acc}/
\identifier{_polyvec_invntt}의
\identifier{verify_matrix_ntt_acc_mode2_cols4_correct}와 actual from-CRT/freeze의
\identifier{verify_crt_freeze_mode2_word_exact}가 없다. challenge leaf
\identifier{verify_tail_m23_highbits_lsb_sampleinball_correct}도 독립적으로 남는다.
이것이 frozen \FrozenPaperGate 다. post-freeze의 pure
\theoremname{verify_mode2_row_product_from_reprs}는 첫 절차 leaf의 대수적 목표를
주지만 actual loop refinement를 대신하지 않는다.

\subsection{9단계 진행: post-freeze actual avec/qj 감사}

\theoremname{rq_mul_coeff_foldr_to_bigi}, full-NTT spectral action, generic row
product, native \(R_q\) 3항 곱과 HAETAE dot representation은 이제
\statusproved 다. actual sampler 쪽도 nonce 515/516, SHAKE128 framing,
little-endian parsing, \(<64513\) rejection과 unchanged storage를 paper \(a\)로
해석하며 paper \(qj\)를 snapshot equation에 별도로 합성한다
\cite{post-freeze-kg-rq-haetae-report,
post-freeze-kg-actual-avec-qj-report}.

현재 첫 목표는 \claimid{OBL-KG-SECURITY-EXPANDVECA-SEMANTICS}다. security
\identifier{public_rounding_vector_seed}가 actual expansion과 동일한 semantics를
갖도록 specification을 교체·정제하고, coefficientwise actual-array
representation을 증명해야 한다. 이후에도 paper \(qj\) correction을 별도로
보존하는 actual-carrier width-6 coefficient-product packaging은 최신
\sourcepath{KgFaithfulAugmentedModel.ec}에서 닫혔다. 이 파일은 actual
matrix/secret 구체화, generated HAETAE \identifier{poly_dot} 운송과
\theoremname{actual_faithful_augmented_productE}를 거쳐 actual checked parent의
product \(\equiv qj\pmod{2q}\) 식을 standalone fresh-compiled theorem으로
합성한다. 이 product는 actual matrix/secret만을 입력으로 받고 중간
2-by-3 block을 그 열에서 복원하므로 \FaithfulAugmentedVerdict 로 판정한다.
남은 의무는 이 전용 mod-\(2q\) block semantics를 기존 mod-\(q\), width-4
security model로 잘못 동일시하지 않은 채, synthetic security object를 actual
ExpandVecA semantics로 교체·정제하는 별도 대응 정리다.

각 부모를 닫을 때 자식
정리(child theorem)의 전제가 실제 앞 절차의 결론에서 나오는지 확인해야 한다.
external precondition으로 원하는 등식을 다시 요구하면 합성이 아니라 조건부
재진술이다.

\subsection{동결된 product/replay 선은 별도로 다룬다}

codec 선과 독립적으로 \claimid{OBL-MU-PRODUCT-REPLAY}는 explicit paper
boundary로 남아 있다. 재개하려면 wrapper 추가가 아니라 다음 중 하나가 필요하다.

\begin{itemize}
  \item hash 반환값을 ghost observation에 저장하는 시점까지 보존하는 sound
        relational sequencing lemma;
  \item 두 continuation을 함께 처리하는 product program과 그 semantics;
  \item Sign continuation의 losslessness/termination을 먼저 증명한 뒤 사용할 수
        있는 정당한 one-sided elimination;
  \item hash-time snapshot을 trace가 보존하도록 하는 최소 instrumentation 변경과
        그 actual-to-trace exactness 재증명.
\end{itemize}

어느 선택도 단순한 byte lemma가 아니다. program logic 설계와 theorem surface
trade-off를 먼저 검토해야 한다.

\subsection{순수 대수학자가 바로 기여할 수 있는 문제}

\begin{question}[rANS를 제한된 전단사(restricted bijection)로 패키징하기]
집합(set)
\[
  \mathcal X_s=\{x:2^{23}\le x<2^{31}\},\qquad
  \mathcal N_s=\{r:1\le r<x_{\max}(s)\}
\]
와 0/1/2-byte 섬유(fiber)를 사용해 normalization--fast-step을 명시적인
부분 전단사(partial bijection)로 패키징할 수 있는가? 현재 개별 lemma를 actual
loop invariant에서 재사용하기 쉬운 하나의 귀납 정리(induction theorem)로 묶는
것이 목표다.
\end{question}

\begin{question}[표현의 몫(quotient)과 정준 단면(canonical section)]
signature 각 구성요소에 대해 raw word 공간(word space), 바이트 몫(byte quotient),
정준 대표자계(canonical system of representatives), decoder 단면(section)을 한
번에 기술하는 공통 구조를 만들 수 있는가?
새 추상화는 실제 기존 술어를 재사용하고, prefix/HBZ/향후 \(h\) codec의
중복만 제거해야 한다.
\end{question}

\begin{question}[ExpandVecA, \(R_q\), mod-\(2q\) 표현]
닫힌 NTT·\(R_q\)·HAETAE 결과의 정확한 명제(proposition)/전제(premise)를
유지하면서 actual SHAKE128/rejection \identifier{avec}를 faithful security
matrix로 어떻게 표현할 수 있는가? paper \(qj\) correction을 modulo-\(q\)
vector와 혼합하지 않고, 새 actual-carrier structured coefficient-product를
augmented matrix/secret만의 generic ring \identifier{matrix_vec_mul} 및 security
model에 어떻게 연결할 수 있는가? Verify 쪽은
동일한 순수 4열 행곱을 actual NTT/CRT/freeze procedure theorem에 연결하는 별도
가환사각형(commutative square)이 필요하다.
\end{question}

\begin{question}[high/low decomposition의 정수 경계(integer bound)]
논문의 \(r=\alpha\operatorname{HighBits}(r)+\operatorname{LowBits}(r)\)와
실제 고정폭 rounding/sign-extension/packing 사이의
정준 대표자(canonical representative) 조건을 정확히
분리할 수 있는가? 이 작업은 전체 challenge input과 norm/hint correctness의
대수적 중심이다.
\end{question}

\begin{question}[rejection distribution]
partial map으로 모델링된 retry loop를 확률 커널(probability kernel)로 올리고,
수락 분포(accepted distribution)와 termination/loss bound를 논문 sampler에
연결할 수 있는가?
결정론적 Hoare 정리(deterministic Hoare theorem)와
분포 정리(distribution theorem)를 섞지 않는 interface 설계가 필요하다.
\end{question}

\subsection{대수기하학자의 기술이 유용한 곳}

현재 작업이 스킴 이론적(scheme-theoretic) 증명은 아니지만 다음 습관은 직접
유용하다.

\begin{itemize}
  \item 전체 객체(object) 대신 필요한 제한(restriction)만 비교하는
        국소--대역(local-to-global) 사고;
  \item 서로 다른 표현 좌표(chart) 사이의 전이사상(transition map)과
        호환성(compatibility) 도식을 명시하는 습관;
  \item 몫(quotient)과 대표 단면(representative section)을 분리하는 습관;
  \item 섬유(fiber)가 비어 있지 않음을 witness로 확인하고,
        조건부 정리(conditional theorem)의
        non-vacuity를 점검하는 습관;
  \item 큰 목표를 가환사각형(commutative square)과 장애물(obstruction)로
        분해하는 습관.
\end{itemize}

단, 이 습관을 ``현재 코드가 층(sheaf)을 형식화한다''는 주장으로 바꾸면 안 된다.
실제 proof term은 배열, word, 메모리, procedure semantics 위에 있다.

\subsection{새 정리(theorem)를 제출할 때의 판정 체크리스트}

\begin{enumerate}
  \item 정리(theorem)의 actual/generated/pure 경계를 이름과 주석에 명시했는가?
  \item 부모 결론을 전제로 다시 가정하지 않았는가?
  \item termination, partial correctness, losslessness를 구분했는가?
  \item canonicality, address range, no-wrap, separation 전제가 모두 보이는가?
  \item 반환값, global state, ghost observation 중 무엇을 비교하는지 명시했는가?
  \item 사용하지 않는 배열·메모리 구간의 frame이 필요한가?
  \item 전제의 concrete witness 또는 실제 predecessor가 있어 non-vacuous한가?
  \item authored axiom, admit/sorry/abort, debug declaration이 없는가?
  \item manifest에 들어가 fresh \identifier{-no-eco} compile되는가?
  \item claim ledger와 theorem graph의 부모 상태를 과장 없이 갱신했는가?
\end{enumerate}

이 체크리스트를 통과해야 ``수학적으로 그럴듯함''이 ``현재 구현에 대해
검증됨''으로 바뀐다.
